\scanpage{76}
\mspage{72}
\uline{10 juillet}

Je reprends \unsure{du début|dès début} le formalisme des disjoncteurs, spatios.

1. \uline{\unsure{Disjoncteurs|Disjoncteur}}, \inserted{\uline{\unsure{supports|support}}.} : Ensemble $E$ muni d'une relation $\mathrel{|\unsure{\circ|0}|}$ \inserted{(dite \uline{rel. de disjonction})}, on \deleted{symétrique. \unsure{O}} \inserted{suppose}

\[
\begin{aligned}
&\text{\uline{Disj.} 1}\qquad x\mathrel{|\circ|}y\ \unsure{\longrightarrow|\text{est}}\ \text{symétrique}\\
&\text{\uline{Disj.} 2}\qquad x\mathrel{|\circ|}x\Longrightarrow x\mathrel{|\circ|}y\quad\forall y\in E
\end{aligned}
\]

% slanted
\inserted{\uline{Structure de disjonction}.} \inserted{On désigne par $E_{0}$ (``\deleted{\unsure{singulariseur|singulariser}}\unsure{singularité|défectuosité}'' du disj.) l'ens. \inserted{des $x$}}
Disjoncteur antiréflexif : quand $\mathrel{|\circ|}$ est antiréflexif (ce qui implique Disj. 2), i.e. $E_{0}=\emptyset$.

\uline{Relation à notion de graphe} (strict). On associe à une str. de disj. $\mathrel{|\circ|}$ sur $E$, une structure de graphe, en prenant comme ens. d'arêtes
\[
A\subset\unsure{\mathcal{P}|\mathrm{P}}_{2}(E)
\]
l'ens. suivant
\[
A=\bigl\{\{x,y\}\in\mathcal{P}_{2}(E)\mid x\mathrel{\overline{|\circ|}}y\bigr\}.
\]

On récupère \inserted{presque} la structure de disjoncteur à partir de la structure de graphe, \deleted{\unsure{en partant}} \inserted{il faut de} plus se donner la partie $E_{0}$ de l'ens.
\[
\deleted{x\mathrel{|\circ|}y\Longleftrightarrow\bigl(\{x,y\}\in A,\ \text{ou }x=y\text{ et }x\ \text{\uline{sommet}}}
\]
des \deleted{\unsure{points}} \inserted{sommets} isolés du graphe, \deleted{\unsure{pour} \unsure{élément} $E_{0}$ \unsure{isolé du graphe} $(E,A)$\unsure{)}}.

\uline{La structure de disjoncteur antiréflexif}$^{\ast}$ \uline{équivaut à celle de graphe strict}
% slanted
\margin{$^{\ast}$ C'est le cas le plus important, bien sûr.}

(On trouve ainsi une corr. bi-univ. entre structures de disjonction sur $E$, et structures de graphe strict, i.e. parties de $\mathcal{P}_{2}(E)$.

La structure de disjoncteur donne un \unsure{formalisme} de \inserted{$\operatorname{cosupp}^{\circ}$ ou} $\operatorname{cosupp}_{E}^{\circ}$ (ou $\operatorname{cosupp}_{\underline{E}}^{\circ}$, \unsure{où} $\deleted{\unsure{E}}\underline{E}=(E,\mathrel{|\circ|}_{E})$) et $\operatorname{supp}_{E}^{\circ}=\operatorname{cosupp}_{E}^{\circ}\circ\operatorname{cosupp}_{E}^{\circ}$ \deleted{\unsure{et}} de $\mathcal{P}(E)$, \unsure{d'est|dont|d'où|c'est}
% archive sheet number: 75
